\documentclass[11pt]{article}

\usepackage[margin=1in]{geometry}
\usepackage{amsmath,amssymb}
\usepackage{booktabs}
\usepackage{array}
\usepackage{listings}
\usepackage{xcolor}
\usepackage{hyperref}
\hypersetup{colorlinks=true,linkcolor=blue!50!black,urlcolor=blue!50!black}

\lstset{
  basicstyle=\ttfamily\small,
  breaklines=true,
  columns=fullflexible,
  keywordstyle=\color{blue!60!black},
  commentstyle=\color{green!40!black},
  showstringspaces=false,
  language=Python,
  frame=single,
  framesep=4pt,
  rulecolor=\color{black!25},
}

\newcommand{\hn}{\hat{n}}
\newcommand{\hx}{\hat{x}}
\newcommand{\hg}{\hat{g}}
\newcommand{\hS}{\hat{S}}
\newcommand{\hV}{\hat{V}}
\newcommand{\hD}{\hat{D}}
\newcommand{\hT}{\hat{T}}
\newcommand{\hf}{\hat{f}}
\newcommand{\code}[1]{\texttt{#1}}

\title{The \code{scipy} Collocation Solver in \code{src/solver\_nondim.py}\\
       \large A flux-variable BVP formulation of the non-dimensional
       Saarelma--Connor pedestal model}
\author{}
\date{}

\begin{document}
\maketitle

\begin{abstract}
\code{src/solver\_nondim.py} solves the coupled three-equation
Saarelma--Connor pedestal model in the non-dimensional form of
Appendix~A.8 of the writeup.  It offers two independent discretisations,
selected by the \code{solver\_structure} keyword of
\code{solve\_coupled\_nondim}: a Firedrake finite-element path
(\code{"firedrake"}, the default) and a SciPy collocation path
(\code{"scipy"}), implemented in
\code{solve\_coupled\_nondim\_scipy}.  This note documents the SciPy
path: the change of variables that turns the mixed second-order system
into an explicit first-order system suitable for
\code{scipy.integrate.solve\_bvp}, the placement of the four boundary
conditions, the Picard loop on the frozen KBM diffusivity, and the
construction of the initial guess.
\end{abstract}

\tableofcontents

\section{Overview}

The physical model is a coupled system for the electron density $n_e$ and
two neutral populations --- Franck--Condon ($n_{\mathrm{FC}}$) and
charge-exchange ($n_{\mathrm{CX}}$) neutrals --- on the pedestal interval
$x \in [x_{\mathrm{inner}}, 0]$, with $x = 0$ the separatrix.  The class
\code{saarelma\_connor\_nondim} subclasses \code{saarelma\_connor} and
inherits all equilibrium, kinetic and cross-section setup unchanged; only
the solve is re-implemented.

The SciPy path performs the following steps.

\begin{enumerate}
  \item Build the equilibrium/kinetic coefficient grids and the ETG
        coefficient (\code{\_ensure\_firedrake\_coefficient\_grids},
        \code{construct\_C\_ETG}) --- pure NumPy, no Firedrake.
  \item Resolve the boundary-condition data
        (\code{bc\_ig\_helpers.resolve\_ne\_bcs}).
  \item Compute the four reference scales and cache them
        (\code{\_set\_nondim\_scales}), Sec.~\ref{sec:nondim}.
  \item Wrap every frozen coefficient in a linear interpolant of $\hx$
        (\code{\_scipy\_coefficients\_nondim}), Sec.~\ref{sec:coeff}.
  \item Change variables from $(\hn_e, \hn_{\mathrm{FC}},
        \hn_{\mathrm{CX}})$ to the flux state
        $Y = (\hn_e, \Phi, U, W)$, Sec.~\ref{sec:flux}.
  \item Build the electron and neutral initial guesses with the helpers
        shared with the Firedrake path, then convert the neutrals to
        fluxes, Sec.~\ref{sec:guess}.
  \item Iterate \code{solve\_bvp} inside a Picard loop that refreezes the
        KBM diffusivity, Sec.~\ref{sec:picard}.
  \item Recover the neutral densities from the fluxes and restore SI
        units, Sec.~\ref{sec:output}.
\end{enumerate}

\section{Non-dimensionalisation}
\label{sec:nondim}

Four free scales are chosen from separatrix values
(\code{\_set\_nondim\_scales}):
\begin{align}
  L    &= |x_{\mathrm{inner}}|                &&\text{(m)},   \\
  n_0  &= n_e(0)                              &&(\mathrm{m^{-3}}), \\
  T_0  &= \bigl(T_e(0) + T_i(0)\bigr)\,e      &&\text{(J)},   \\
  S_0  &= S_i(0)                              &&(\mathrm{m^3\,s^{-1}}),
\end{align}
with the derived scales
\begin{equation}
  \tau = \frac{1}{n_0 S_0}, \qquad
  [V]_0 = \frac{L}{\tau} = L n_0 S_0, \qquad
  [D]_0 = \frac{L^2}{\tau} = L^2 n_0 S_0, \qquad
  [p]_0 = n_0 T_0 .
\end{equation}
Hatted variables are $\hx = x/L \in [-1,0]$, $\hn_e = n_e/n_0$,
$\hn_{\mathrm{FC}} = \langle n_{\mathrm{FC}}\rangle/n_0$,
$\hn_{\mathrm{CX}} = \langle n_{\mathrm{CX}}\rangle/n_0$,
$\hS_i = S_i/S_0$, $\hS_{\mathrm{CX}} = S_{\mathrm{CX}}/S_0$,
$\hV = |V|/[V]_0$, and $\hg = \langle |\nabla r|^2 \rangle$ (already
dimensionless).  The transport coefficients scale as
\begin{equation}
  \hat{C}_{\mathrm{ETG}} = \frac{C_{\mathrm{ETG}}}{L^2 n_0^2 S_0},\qquad
  \hD_{\mathrm{NEO}}     = \frac{D_{\mathrm{NEO}}}{[D]_0},\qquad
  \hat{A}_{\mathrm{KBM}} = \frac{A_{\mathrm{KBM}}}{[D]_0},\qquad
  \hat{B}_{\mathrm{KBM}} = \frac{B_{\mathrm{KBM}} T_0}{L^3 S_0}.
\end{equation}
By construction $\hn_e(0) = 1$ and $\hS_i(0) = 1$.

\section{Governing equations}

In strong (divergence) form the non-dimensional system is
\begin{align}
  \frac{d}{d\hx}\Bigl[ \hf\, \frac{d\hn_e}{d\hx} \Bigr]
    &= -\,\hn_e\,\hS_i\,\bigl(\hn_{\mathrm{FC}} + \hn_{\mathrm{CX}}\bigr),
    \label{eq:ne} \\[2pt]
  \hV_{\mathrm{FC}}\,\frac{d}{d\hx}
    \bigl[ f_{\mathrm{FC}}\,\hg\,\hn_{\mathrm{FC}} \bigr]
    &= \hn_e\,\bigl(\hS_i + \hS_{\mathrm{CX}}\bigr)\,\hn_{\mathrm{FC}},
    \label{eq:nfc} \\[2pt]
  \frac{d}{d\hx}
    \bigl[ \hV_{\mathrm{CX}}\,f_{\mathrm{CX}}\,\hg\,\hn_{\mathrm{CX}} \bigr]
    &= \hn_e\Bigl(\hS_i\,\hn_{\mathrm{CX}}
        - \tfrac{1}{2}\,\hS_{\mathrm{CX}}\,\hn_{\mathrm{FC}}\Bigr),
    \label{eq:ncx}
\end{align}
which are the strong counterparts of the weak residuals
\code{F1\_hat}, \code{F2\_hat}, \code{F3\_hat} in the module docstring.
The \emph{conductance} $\hf$ collects all the diffusive channels:
\begin{equation}
  \hf(\hx, \hn_e)
    = \hg\,\bigl(\hD_{\mathrm{NEO}} + \hD_{\mathrm{KBM}}\bigr)
    + \hg\,\frac{\hat{C}_{\mathrm{ETG}}}{\hn_e}.
  \label{eq:conductance}
\end{equation}
The $\hat{C}_{\mathrm{ETG}}/\hn_e$ term makes $\hf$ density-dependent;
this is the only nonlinearity in the $\hn_e$ flux, because the SciPy
path runs with $\hat{B}_{\mathrm{KBM}} = 0$
(Sec.~\ref{sec:picard}), so no $(\hn_e')^2$ term survives.

Equations \eqref{eq:ne}--\eqref{eq:ncx} are second order in $\hn_e$ and
first order in each neutral: a fourth-order system requiring four
boundary conditions, exactly as in Saarelma et al.~(2023), Sec.~2.3.

\section{The flux change of variables}
\label{sec:flux}

\code{solve\_bvp} requires an explicit first-order system
$\mathbf{y}' = \mathbf{F}(\hx, \mathbf{y})$.  The naive reduction
$y = (\hn_e, \hn_e', \dots)$ would require expanding
$d/d\hx[\hf \hn_e']= \hf \hn_e'' + \hf' \hn_e'$, i.e.\ a coefficient
derivative $d\hf/d\hx$ evaluated off the coarse p-file grid --- the
accuracy-limiting step noted in \code{solver\_sc}'s module docstring.
The SciPy path avoids it entirely by carrying the three \emph{fluxes} as
state variables:
\begin{equation}
  Y = \bigl[\, \hn_e,\ \Phi,\ U,\ W \,\bigr],
  \qquad
  \begin{aligned}
    \Phi &= \hf\,\frac{d\hn_e}{d\hx}
          &&\text{(electron flux)},\\
    U    &= f_{\mathrm{FC}}\,\hg\,\hn_{\mathrm{FC}}
          &&\text{(FC flux}/\hV_{\mathrm{FC}}\text{)},\\
    W    &= \hV_{\mathrm{CX}}\,f_{\mathrm{CX}}\,\hg\,\hn_{\mathrm{CX}}
          &&\text{(CX flux)}.
  \end{aligned}
\end{equation}
Substituting into \eqref{eq:ne}--\eqref{eq:ncx} and using the inverse
relations
\begin{equation}
  \hn_{\mathrm{FC}} = \frac{U}{f_{\mathrm{FC}}\,\hg},
  \qquad
  \hn_{\mathrm{CX}} = \frac{W}{\hV_{\mathrm{CX}}\,f_{\mathrm{CX}}\,\hg},
  \label{eq:recover}
\end{equation}
gives the first-order system integrated by the solver:
\begin{equation}
\left\{
\begin{aligned}
  \frac{d\hn_e}{d\hx} &= \frac{\Phi}{\hf(\hx,\hn_e)},
    \\
  \frac{d\Phi}{d\hx}  &= -\,\hn_e\,\hS_i\,
    \bigl(\hn_{\mathrm{FC}} + \hn_{\mathrm{CX}}\bigr),
    \\
  \frac{dU}{d\hx}     &= \frac{\hn_e\,
    \bigl(\hS_i + \hS_{\mathrm{CX}}\bigr)\,\hn_{\mathrm{FC}}}
    {\hV_{\mathrm{FC}}},
    \\
  \frac{dW}{d\hx}     &= \hn_e\Bigl(\hS_i\,\hn_{\mathrm{CX}}
    - \tfrac{1}{2}\,\hS_{\mathrm{CX}}\,\hn_{\mathrm{FC}}\Bigr).
    
\end{aligned}
\right.
\label{eq:odesys}
\end{equation}
\textbf{No coefficient derivative appears anywhere.}  Every coefficient
enters algebraically, so the divergence form of the equations is kept
intact --- the same property the Firedrake weak form enjoys.  This is the
central design point of the SciPy path.

In the code this is the closure \code{ode(xq, Y)}, which returns
the four right-hand sides of \eqref{eq:odesys} stacked with \code{np.vstack}.  Note
that \code{ode} is written to be vectorised over the collocation nodes
\code{xq}, as \code{solve\_bvp} requires.

\section{Frozen coefficients}
\label{sec:coeff}

All frozen coefficients live on the SI grid \code{self.x\_init}.
\code{\_scipy\_coefficients\_nondim} returns a dictionary of callables of
$\hx$, each of which internally applies $x = L\hx$ and multiplies by the
appropriate inverse scale:

\begin{center}
\begin{tabular}{@{}llll@{}}
\toprule
Key & Symbol & Source array & Scaling \\
\midrule
\code{g}      & $\hg$                      & \code{gradr2\_fsa}       & $1$ (dimensionless) \\
\code{Si}     & $\hS_i$                    & \code{S\_i\_pres}        & $1/S_0$ \\
\code{Scx}    & $\hS_{\mathrm{CX}}$        & \code{S\_cx\_pres}       & $1/S_0$ \\
\code{Vcx}    & $\hV_{\mathrm{CX}}$        & $|$\code{V\_cx\_pres}$|$ & $1/[V]_0$ \\
\code{fFC}    & $f_{\mathrm{FC}}$          & \code{fFC}               & $1$ \\
\code{fCX}    & $f_{\mathrm{CX}}$          & \code{fCX}               & $1$ \\
\code{D\_NEO} & $\hD_{\mathrm{NEO}}$       & \code{D\_NEO}            & $1/[D]_0$ \\
\code{C\_ETG} & $\hat{C}_{\mathrm{ETG}}$   & \code{C\_ETG}            & $1/(L^2 n_0^2 S_0)$ \\
\code{VFC}    & $\hV_{\mathrm{FC}}$        & $|$\code{V\_FC}$|$ (scalar) & $1/[V]_0$ \\
\bottomrule
\end{tabular}
\end{center}

Interpolation is \emph{linear} with linear extrapolation
(\code{interp1d(..., kind='linear', fill\_value='extrapolate')}),
deliberately matching the \code{np.interp} sampling the Firedrake path
uses to fill its coefficient \code{Function}s.  Both discretisations
therefore see the identical coefficient reconstruction, which makes a
solver-to-solver comparison meaningful.

The KBM diffusivity $\hD_{\mathrm{KBM}}$ is \emph{not} in this
dictionary: it is refrozen every Picard iteration and enters through a
separate interpolant \code{D\_KBM\_x} rebuilt inside the loop.

\section{Boundary conditions}
\label{sec:bcs}

Four residuals are required.  Three are always the same Dirichlet
conditions at the separatrix $\hx = 0$, expressed in flux variables via
\eqref{eq:recover}:
\begin{align}
  \hn_e(0) &= \hn_e^{(0)} \;=\; n_e(0)/n_0 \;=\; 1, \\
  U(0)     &= U_0 \;=\; \hn_{\mathrm{FC}}^{(0)}\,
                 f_{\mathrm{FC}}(0)\,\hg(0), \\
  W(0)     &= W_0 \;=\; \hn_{\mathrm{CX}}^{(0)}\,
                 \hV_{\mathrm{CX}}(0)\,f_{\mathrm{CX}}(0)\,\hg(0).
\end{align}
The fourth is the prescribed density gradient, converted to a flux
condition through $\Phi = \hf\,\hn_e'$.  Its placement is controlled by
\code{ne\_grad\_bc\_loc}:
\begin{itemize}
  \item \code{"inner"} (default): imposed at $\hx = -1$,
        \[
          \Phi(-1) = \hf\bigl(-1,\ \hn_e(-1)\bigr)\;
          \widehat{\left.\frac{dn_e}{dx}\right|}_{\mathrm{inner}} ,
        \]
        one condition at each end of the domain;
  \item \code{"outer"}: imposed at $\hx = 0$ alongside the Dirichlet
        value, $\Phi(0) = \hf(0, \hn_e(0))\,\widehat{dn_e/dx}|_{0}$,
        leaving the inner boundary completely free.  This is the
        boundary-condition set of Saarelma et al.~(2023), Sec.~2.3.
\end{itemize}
The dimensionless gradient is $\widehat{dn_e/dx} = L\,(dn_e/dx)/n_0$
(\code{hat\_dne\_dx\_bc}).

\paragraph{Why the SciPy path needs no shooting.}
This is the practical advantage of collocation here.  In the Galerkin
discretisation, a strong Dirichlet condition at $\hx = 0$ forces
$v_e(0) = 0$, so a boundary flux term at the separatrix vanishes
identically and both conditions cannot be imposed directly; the Firedrake
path therefore recasts the inner flux as an unknown Lagrange multiplier
and runs a secant iteration on it (\code{\_shoot\_outer\_grad\_nondim}).
\code{solve\_bvp} only cares about the correct \emph{number} of residuals
and is indifferent to which end each is evaluated at, so
\code{ne\_grad\_bc\_loc="outer"} costs nothing extra --- it is one line in
the \code{bc(Ya, Yb)} closure.  In that closure \code{Ya} is the state at
$\hx=-1$ and \code{Yb} the state at $\hx=0$.

\section{Initial guess and neutral seeds}
\label{sec:guess}

The collocation grid is uniform and ascending,
$\hx \in \{-1, \dots, 0\}$ with \code{x\_res} points.  Both the density
and the neutral guesses come from the shared helpers in
\code{src/bc\_ig\_helpers.py}, which are the \emph{same} routines the
Firedrake path calls, on the same SI grid.  That is deliberate: the
module docstring of \code{bc\_ig\_helpers} states that every solver in
the repository resolves its boundary conditions and builds its initial
guesses there so the discretisations cannot drift apart.  The two
paths therefore start from identical SI profiles and accept identical
options.

\subsection{Electron density}

\code{build\_ne\_initial\_guess(model, x\_grid, initial\_guess, bcs,
tanh\_width, tanh\_center)} returns $n_e^{\mathrm{init}}$ in SI on the
requested grid, with \code{initial\_guess} one of \code{"linear"},
\code{"pfile"}, \code{"tanh"} or \code{"manual EPEDNN loop"}.  Every
shape is anchored on the resolved boundary conditions: it meets
\code{bcs.ne\_outer} at the separatrix and rises to
\code{bcs.ne\_inner\_guess} at the pedestal top.  The scipy path then
normalises and floors it,
\begin{equation}
  N_g = \max\bigl(n_e^{\mathrm{init}}/n_0,\ \varepsilon\bigr),
  \qquad \varepsilon = \code{ne\_floor}\ (\text{default } 10^{-8}),
\end{equation}
the floor guarding the $\hat{C}_{\mathrm{ETG}}/\hn_e$ division.

\subsection{Neutrals}
\label{sec:neutralguess}

\code{build\_neutral\_initial\_guess(model, x\_grid, ne\_init, nFC\_ic,
nCX\_ic)} returns the pair $(n_{\mathrm{FC}}^{\mathrm{init}},
n_{\mathrm{CX}}^{\mathrm{init}})$ in SI.  All work is done on the grid
sorted in \emph{descending} $x$ --- integrating inward from the
separatrix, the end that carries the Dirichlet data
$n_{\mathrm{FC}}(0) = $ \code{nFC\_x0} and $n_{\mathrm{CX}}(0) = $
\code{nCX\_x0} --- and the results are scattered back into the caller's
ordering, so an unsorted \code{x\_grid} is handled correctly.

\paragraph{\code{nFC\_ic="solve"}.}
At frozen $n_e$, Eq.~(14) of Saarelma et al.~(2023),
\begin{equation}
  |V_{\mathrm{FC}}|\,\frac{d}{dx}
    \bigl[f_{\mathrm{FC}}\,n_{\mathrm{FC}}\bigr]
  = n_e\,\bigl(S_i + S_{\mathrm{CX}}\bigr)\,n_{\mathrm{FC}},
\end{equation}
is exactly linear and homogeneous, so an integrating factor anchored at
the separatrix gives it in closed form:
\begin{equation}
  n_{\mathrm{FC}}(x) = n_{\mathrm{FC}}(0)\,
    \exp\!\left[\int_0^{x}
      \frac{n_e\,(S_i + S_{\mathrm{CX}})}
           {f_{\mathrm{FC}}\,|V_{\mathrm{FC}}|}\,dx'\right].
  \label{eq:fcguess}
\end{equation}
The integral is evaluated with \code{cumulative\_trapezoid} over the
descending grid, so $dx < 0$ and the guess decays inward --- the
boundary-layer behaviour the collocation mesh has to resolve.  A
negative result raises \code{ValueError}.

\paragraph{\code{nCX\_ic="solve"}.}
The CX equation, Eq.~(10),
\begin{equation}
  |V_{\mathrm{CX}}|\,\frac{d}{dx}
    \bigl[f_{\mathrm{CX}}\,g\,n_{\mathrm{CX}}\bigr]
  = n_e\Bigl(S_i\,n_{\mathrm{CX}}
      - \tfrac{1}{2}S_{\mathrm{CX}}\,n_{\mathrm{FC}}\Bigr),
\end{equation}
is linear but \emph{inhomogeneous}: it is driven by the FC guess just
built.  With $u = f_{\mathrm{CX}}\,g\,n_{\mathrm{CX}}$ and the inward
coordinate $\tau = -x \ge 0$,
\begin{equation}
  \frac{du}{d\tau} = -P(\tau)\,u + Q(\tau),
  \qquad
  P = \frac{n_e S_i}{|V_{\mathrm{CX}}|\,f_{\mathrm{CX}}\,g},
  \qquad
  Q = \frac{n_e\,S_{\mathrm{CX}}\,n_{\mathrm{FC}}}{2\,|V_{\mathrm{CX}}|},
\end{equation}
noting that $f_{\mathrm{CX}}$ does \emph{not} appear in $Q$.  The
integrating factor $\nu(\tau) = \exp\!\int_0^\tau P\,d\tau'$ gives
\begin{equation}
  u(\tau) = \frac{u(0) + \int_0^\tau \nu\,Q\,d\tau'}{\nu(\tau)},
  \qquad
  u(0) = f_{\mathrm{CX}}(0)\,g(0)\,n_{\mathrm{CX}}(0),
  \qquad
  n_{\mathrm{CX}} = \frac{u}{f_{\mathrm{CX}}\,g}.
  \label{eq:cxguess}
\end{equation}
Keeping the $Q$ term matters: the CX population in this model is
\emph{fed} by charge exchange on the FC neutrals, so a purely
homogeneous seed misses its source entirely.

\paragraph{Manual and scaled options.}
\code{nFC\_ic="manual EPEDNN loop"} and
\code{nCX\_ic="manual EPEDNN loop"} interpolate the tabulated profiles
\code{model.nFC\_manual} / \code{model.nCX\_manual}, which are given on
\code{model.psi\_N\_n\_manual}, onto $x$ via the
$\psi_N \mapsto x$ map \code{model.psi\_N\_pres} $\to$
\code{model.x\_init}.  These are the profiles handed back by an outer
EPED-NN iteration, so a coupled solve can be restarted from the
previous outer step rather than from an analytic guess.
\code{nCX\_ic="scale nFC"} simply takes
$n_{\mathrm{CX}}^{\mathrm{init}} = n_{\mathrm{FC}}^{\mathrm{init}}\,
n_{\mathrm{CX}}(0)/n_{\mathrm{FC}}(0)$.
The two choices are independent: any \code{nFC\_ic} may be combined with
any \code{nCX\_ic}, including a manual FC guess driving a solved CX
guess through $Q$ in \eqref{eq:cxguess}.

\subsection{Conversion to the flux state}

The helper returns SI \emph{densities}, but the collocation system
carries fluxes, so the seeds are converted through \eqref{eq:recover}:
\begin{equation}
  U_g = \frac{n_{\mathrm{FC}}^{\mathrm{init}}}{n_0}\,
        f_{\mathrm{FC}}\,\hg,
  \qquad
  W_g = \frac{n_{\mathrm{CX}}^{\mathrm{init}}}{n_0}\,
        \hV_{\mathrm{CX}}\,f_{\mathrm{CX}}\,\hg,
  \qquad
  \Phi_g = \hf(\hx, N_g)\,\nabla N_g,
\end{equation}
the last with \code{np.gradient}.  Stacking gives the
$4 \times$\code{x\_res} array \code{Y\_guess} that seeds the first
\code{solve\_bvp} call.

\section{The Picard loop on the KBM diffusivity}
\label{sec:picard}

The KBM contribution is handled by outer (Picard) iteration rather than
being folded into the Newton solve, following Saarelma et
al.~(2023) Eqs.~(24)--(25).  Only \code{picard\_gate\_mode="average"} is
supported here.  In that mode
\code{calc\_pressure\_quantities\_nondim} computes the Connor--Hastie
$\alpha$ in SI units from the current density,
\begin{equation}
  p = n_e\,(T_e + T_i)\,e, \qquad
  \alpha = \alpha_{\mathrm{nodp}}\,\frac{dp}{dx},
\end{equation}
gates on the pedestal-averaged value, and returns a single whole-grid
diffusivity
\begin{equation}
  D_{\mathrm{KBM}} =
  \begin{cases}
    (\bar{\alpha} - \alpha_{\mathrm{crit}})\,G_{\mathrm{KBM}},
      & \bar{\alpha} > \alpha_{\mathrm{crit}},\\
    0, & \text{otherwise},
  \end{cases}
  \qquad
  \hD_{\mathrm{KBM}} = D_{\mathrm{KBM}}/[D]_0 ,
\end{equation}
with $G_{\mathrm{KBM}} = C_{\mathrm{KBM}} c_s \rho_s^2 / a$.  Because the
whole frozen diffusivity sits in the $D$-slot with
$\hat{B}_{\mathrm{KBM}} = 0$, the conductance
\eqref{eq:conductance} depends on $\hx$ and $\hn_e$ only --- never on
$\hn_e'$ --- which is precisely what makes the inversion
$\Phi \mapsto \hn_e' = \Phi/\hf$ in the first line of \eqref{eq:odesys} explicit.

\paragraph{Why \code{"majority"} is rejected.}
The \code{"majority"} gate freezes the local $(A,B)$ KBM structure at
every point, which reintroduces a $\hat{B}_{\mathrm{KBM}}$ term
proportional to $(\hn_e')^2$.  The conductance would then depend on
$\hn_e'$, and recovering $\hn_e'$ from $\Phi$ would become a per-point
root-find inside the ODE right-hand side.  That is not implemented; the
code raises \code{NotImplementedError} and directs the user to
\code{solver\_structure="firedrake"}.

\paragraph{Loop structure.}
Each iteration:
\begin{enumerate}
  \item rebuild the interpolant \code{D\_KBM\_x} from the current
        $\hD_{\mathrm{KBM}}$;
  \item interpolate the previous solution onto the base grid as the
        starting guess (\code{solve\_bvp} adapts its own mesh, so
        \code{sol.x} generally differs from \code{hat\_x});
  \item call \code{solve\_bvp(ode, bc, hat\_x, Y\_start, tol=bvp\_tol,
        max\_nodes=bvp\_max\_nodes)}; a failure raises immediately;
  \item measure the relative density change
        $\delta = \|N^{\mathrm{new}} - N^{\mathrm{old}}\|_\infty /
        \|N^{\mathrm{new}}\|_\infty$;
  \item refreeze $\hD_{\mathrm{KBM}}$ from the new profile, with optional
        under-relaxation
        \[
          \hD_{\mathrm{KBM}}^{\,\mathrm{new}}
            = \omega\,\hD_{\mathrm{KBM}}^{\,\ast}
            + (1-\omega)\,\hD_{\mathrm{KBM}}^{\,\mathrm{old}},
          \qquad \omega = \code{picard\_relax} \in (0,1];
        \]
  \item append a history record and test convergence.
\end{enumerate}
Convergence requires \emph{both} $\delta < $ \code{picard\_rtol}
\emph{and} an unchanged gate state between successive iterations.  The
second condition matters: the gate is a discontinuous function of
$\bar\alpha$, so a run can otherwise oscillate between a KBM-on and a
KBM-off branch with a small $\delta$ on each.  Failure to converge within
\code{picard\_max\_it} raises a \code{RuntimeError} suggesting a larger
iteration budget or $\omega < 1$.  Diagnostics land in
\code{self.picard\_info} and \code{self.kbm\_info}.

\section{Recovery and output}
\label{sec:output}

On the converged adaptive mesh \code{sol.x}, the neutral densities are
recovered algebraically from \eqref{eq:recover} and everything is
rescaled to SI:
\begin{equation}
  x = L\,\hx, \qquad
  n_e = n_0\,\hn_e, \qquad
  n_{\mathrm{FC}} = n_0\,\hn_{\mathrm{FC}}, \qquad
  n_{\mathrm{CX}} = n_0\,\hn_{\mathrm{CX}}, \qquad
  \frac{dn_e}{dx} = \frac{n_0}{L}\,\frac{\Phi}{\hf}.
\end{equation}
The results are stored on the same attributes the parent solver uses
(\code{x\_sol}, \code{ne\_sol}, \code{nFC\_sol}, \code{nCX\_sol},
\code{dne\_dx\_sol}), so downstream plotting and EPED feeds are
unaffected; the dimensionless profiles are kept alongside as
\code{hat\_x\_sol}, \code{hat\_ne\_sol}, \code{hat\_nFC\_sol},
\code{hat\_nCX\_sol}, and the raw \code{solve\_bvp} object as
\code{self.sol}.  The return value is the tuple
\code{(x\_sol, ne\_sol, nFC\_sol, nCX\_sol, T\_e\_pres, psi\_N\_pres)},
identical to the Firedrake path's.

\section{Algorithm summary}

\begin{lstlisting}[language=Python, caption={Structure of
\code{solve\_coupled\_nondim\_scipy}.}]
# --- setup -----------------------------------------------------------
validate(picard_gate_mode == "average", picard_relax in (0, 1])
_ensure_firedrake_coefficient_grids(x_res);  construct_C_ETG()
nebcs = resolve_ne_bcs(...)          # ne(0), dne/dx at the chosen end
_set_nondim_scales()                 # L, n0, T0, S0 -> tau, [V]0, [D]0
hat_x = linspace(-1, 0, x_res)
c     = _scipy_coefficients_nondim() # callables of hat_x

# --- initial guess (shared helpers: src/bc_ig_helpers.py) ------------
ne_init   = build_ne_initial_guess(self, x_si, initial_guess, nebcs, ...)
N_g       = clip(ne_init / n0, ne_floor, None)
hat_D_KBM = calc_pressure_quantities_nondim(N_g, "average", hat_x)
nFC_init, nCX_init = build_neutral_initial_guess(self, x_si, ne_init,
                                                 nFC_ic, nCX_ic)
U_g       = (nFC_init / n0) * fFC * g          # SI densities -> fluxes
W_g       = (nCX_init / n0) * Vcx * fCX * g
Phi_g     = conductance(hat_x, N_g, hat_D_KBM) * gradient(N_g, hat_x)
Y_prev    = vstack([N_g, Phi_g, U_g, W_g])

# --- Picard loop on the frozen KBM diffusivity -----------------------
for it in 1 .. picard_max_it:
    D_KBM_x = interp1d(hat_x, hat_D_KBM)      # enters ode() and bc()
    Y_start = interp(hat_x, x_prev, Y_prev)   # sol.x is adaptive
    sol     = solve_bvp(ode, bc, hat_x, Y_start,
                        tol=bvp_tol, max_nodes=bvp_max_nodes)
    raise if not sol.success
    dn_rel  = relative_change(sol.y[0], Y_prev[0])
    hat_D_KBM = relax(calc_pressure_quantities_nondim(sol.y[0], ...))
    if dn_rel < picard_rtol and gate_unchanged: converged; break

# --- recover and rescale ---------------------------------------------
nFC = U / (fFC * g);   nCX = W / (Vcx * fCX * g)
x_sol, ne_sol, ... = L * hat_x_sol, n0 * N_sol, ...
\end{lstlisting}

\section{Comparison with the Firedrake path}

\begin{center}
\renewcommand{\arraystretch}{1.2}
\begin{tabular}{@{}p{0.24\textwidth}p{0.34\textwidth}p{0.34\textwidth}@{}}
\toprule
& \code{solver\_structure="firedrake"} & \code{solver\_structure="scipy"} \\
\midrule
Discretisation & CG$_p$ finite elements on a mixed space, Newton on the
weak residuals & Collocation (\code{solve\_bvp}) on an adaptive mesh \\
Unknowns & $(\hn_e, \hn_{\mathrm{FC}}, \hn_{\mathrm{CX}})$ &
$(\hn_e, \Phi, U, W)$ --- fluxes \\
Coefficient derivatives & none (weak form) & none (divergence form kept) \\
Mesh & fixed, \code{x\_res} cells & adaptive, up to \code{bvp\_max\_nodes} \\
Error control & mesh resolution & \code{bvp\_tol} residual tolerance \\
\code{ne\_grad\_bc\_loc="outer"} & secant shooting on the inner flux &
one extra residual, no shooting \\
KBM \code{"average"} & supported & supported \\
KBM \code{"majority"} / inline gate & supported & \code{NotImplementedError} \\
n$_e$ initial guess & \multicolumn{2}{l}{shared
\code{build\_ne\_initial\_guess}} \\
Neutral initial guess & \multicolumn{2}{l}{shared
\code{build\_neutral\_initial\_guess} (\code{nFC\_ic}, \code{nCX\_ic})} \\
Analytic neutral option & \code{neutral\_treatment="analytic"} available &
not applicable (neutrals solved with the system) \\
\bottomrule
\end{tabular}
\end{center}

\section{Practical notes}

\begin{itemize}
  \item \textbf{Entry point.}  Reach the SciPy path through
        \code{solve\_coupled\_nondim(solver\_structure="scipy", ...)}, or
        call \code{solve\_coupled\_nondim\_scipy} directly.  Note that
        \code{solve\_coupled\_nondim} checks
        \code{\_FIREDRAKE\_AVAILABLE} \emph{before} dispatching, so
        reaching the SciPy path through the driver still requires
        Firedrake to be importable even though the path itself never uses
        it; calling \code{solve\_coupled\_nondim\_scipy} directly does
        not.
  \item \textbf{\code{x\_res}} only sets the \emph{initial} collocation
        grid.  \code{solve\_bvp} refines adaptively until the residual
        meets \code{bvp\_tol}, so the returned \code{sol.x} is generally
        finer and non-uniform.  The default \code{x\_res=200} for the
        SciPy path is much larger than the Firedrake default
        \code{x\_res=20} because the initial mesh must at least see the
        neutral boundary layer.
  \item \textbf{\code{ne\_floor}} appears in three places: the initial
        guess, the $\hat{C}_{\mathrm{ETG}}/\hn_e$ division inside
        \code{ode}, and the conductance evaluations inside \code{bc}.  It
        is a division guard, not a physical positivity constraint ---
        nothing prevents \code{solve\_bvp} from returning a slightly
        negative $\hn_e$ if the problem is under-resolved.
  \item \textbf{Overflow in \code{nCX\_ic="solve"}.}  The integrating
        factor $\nu = \exp\!\int_0^\tau P\,d\tau'$ in \eqref{eq:cxguess}
        \emph{grows} inward and is not clipped, so on a strongly
        neutral-opaque case both $\nu$ and $\int \nu Q$ can overflow to
        \code{inf} and their ratio return \code{nan}.  The FC guess
        \eqref{eq:fcguess} has the opposite sign and merely underflows
        to zero, harmlessly.  This is a property of the shared helper
        and so affects both discretisations equally; the mathematically
        equivalent but numerically stable rearrangement is
        $u(\tau) = u_0 e^{-I(\tau)}
         + \int_0^\tau e^{-(I(\tau) - I(s))} Q(s)\,ds$
        with $I(\tau) = \int_0^\tau P$.  If a solve returns \code{nan}
        neutrals, this is the first place to look.
  \item \textbf{Failure modes.}  A \code{solve\_bvp} failure raises
        immediately with \code{sol.message}; the usual causes are too
        coarse an initial mesh, a badly scaled initial guess, or a
        \code{max\_nodes} budget exhausted by an unresolved neutral
        layer.  Increasing \code{x\_res} or loosening \code{bvp\_tol} is
        the first thing to try.
\end{itemize}

\end{document}
